\section{Selection rules for a source}\label{sec:selection}
A source for the localized form is a complex-linear $\Phi_L$ from
$C_c^\infty(I_L)$ into a positive Hilbert space with
$\norm{\Phi_L(f)}^2=Q_L[f]$. The following are necessary conditions on any such
$\Phi_L$, ordered by what they cost to apply. Conditions 1, 2 and 5 are results
of the companion manuscript and are stated here for completeness; conditions 3,
4, 6, 7 and 9 are results of this paper or of the investigation accompanying it.
The point of collecting them is that testing a candidate by matching the terms of
\eqref{eq:target} one at a time is not a weaker version of the right question:
every term can be matched exactly while no source exists.

\begin{enumerate}
\item\textbf{Reflection.} The pairing commutes with $x\mapsto-x$: every source
carries a unitary $J$ with $J\Phi_L(f)=\Phi_L(f(-\cdot))$ and $J^2=\mathrm{id}$,
whose eigenspaces carry the two directions of \eqref{eq:parity-split}.
\item\textbf{Unboundedness.} $\Phi_L$ is unbounded relative to the input
$L^2$ norm, because $Q_L[\chi e^{iNx}]$ grows like $\norm\chi^2\log N$. This
excludes every source with a bounded spectral variable, and in particular the
conformal Wilson pairing, whose operator has spectrum in $[-2,2]$, so that its
kernel is entire in the separation and has no atoms. The criterion is unbounded
spectrum, not open versus closed geometry: the test functions carry no
boundary conditions.
\item\textbf{The spectral weight is the smooth zero density}, to four orders:
the leading logarithm with unit coefficient, the intercept $-\log2\pi$, the
absence of an $|\tau|^{-1}$ term, and then $-1/24$ and $-7/960$, by
\eqref{eq:graded}. A candidate with the right slope and the wrong intercept has
the wrong density of states.
\item\textbf{The archimedean channel spectrum is $\{2n+\tfrac12\}$ with unit
residues}, exactly, by \eqref{eq:gamma-kernel}. Reading $\gammak$ as a channel
spectrum assumes an identification of the arithmetic coordinate with Euclidean
time which a model must supply; where it applies, it says the transverse tower
has spacing $2$ and the quarter shift.
\item\textbf{Every channel decomposition passes the interference bound}
$t\,A[f]+t^{-1}B[f]+C[f]\geq0$ for all $t>0$. Applied to the archimedean energy
and the summed prime references this fails at the indicator of $I_{5/4}$ and by
a growing factor in $L$. Apply it first: it is cheap and it excludes most
proposals.
\item\textbf{The archimedean part is not a lifted circle weight.} A Bloch lift
of a weight on $S^1$ has multiplier $\tau\mapsto w(e^{id\tau})$: continuous,
positive, periodic and bounded, with purely atomic off-diagonal kernel. The
archimedean symbol is unbounded with a smooth, nowhere vanishing off-diagonal
kernel. No finite sum of such lifts, at any parameters, any number of primes and
any quantization parameter, is the archimedean form.
\item\textbf{The contraction is critical, with a margin collapsing in $L$.}
Weil positivity on $I_L$ is a domination whose saturation
$1-\min\operatorname{spec}(Q_L;K_+)$ is below $10^{-9}$ at $L=1$ and below
$10^{-24}$ at $L=3$. Any hypothesis supplying a margin uniform in $L$ ---
a self-map into a fixed compact subset of a cone's interior, a uniform
contraction ratio, a spectral gap --- proves something false.
\item\textbf{The complete identity at $L=1$}, the only sufficient test on this
list, and the expensive one.
\item\textbf{The source is a jump process with L\'evy measure $\mu_L$.} By
Proposition~\ref{prop:jump} the entire non-pole part of the target is the
Dirichlet form of \eqref{eq:levy-measure}: continuous part $2\gammak$, and an
atom of mass $2(\log p)p^{-m/2}$ at each $m\log p$. This subsumes conditions 3
and 4 and replaces them by one, and it removes the appearance of two mechanisms:
the archimedean and prime halves are the continuous and atomic parts of one
measure.
\end{enumerate}

Conditions considered and discarded are worth naming too. A ghost or indefinite
sector for the pole term is not required and cannot exist in a source compatible
in $L$, which is unitarily the evaluation map on the zeros in a positive
$\ell^2$; what survives of that discussion is condition 1. The prime-free
archimedean inequality is implied by $Q_L\geq0$ and so cannot fail unless the
Riemann hypothesis does, and discriminates nothing.
